\documentclass[12pt,a4paper]{article}

\usepackage[utf8]{inputenc}
\usepackage[T1]{fontenc}
\usepackage[brazil]{babel}
\usepackage[margin=2.5cm]{geometry}
\usepackage{graphicx}
\usepackage{booktabs}
\usepackage{longtable}
\usepackage{hyperref}
\usepackage{enumitem}

\title{Projeto 04: Receita Fácil\\Aplicativo de Receitas Culinárias}
\author{Nome do Aluno -- Matrícula: XXXX}
\date{\today}

\begin{document}

\maketitle
\tableofcontents
\newpage

\section{Levantamento de Requisitos}

O Receita Fácil é um aplicativo mobile destinado a usuários que desejam buscar receitas culinárias, consultar ingredientes e visualizar o modo de preparo de pratos.

A técnica considerada para o levantamento inicial foi a observação de necessidades comuns de usuários que procuram receitas pelo celular.

\subsection{Requisitos Funcionais}

\begin{longtable}{p{2cm}p{10cm}}
\toprule
Código & Descrição\\
\midrule
RF-01 & O sistema deve permitir buscar receitas pelo nome.\\
RF-02 & O sistema deve apresentar imagem e nome das receitas.\\
RF-03 & O sistema deve apresentar os ingredientes da receita.\\
RF-04 & O sistema deve apresentar o modo de preparo.\\
RF-05 & O sistema deve permitir favoritar receitas.\\
RF-06 & O sistema deve armazenar os favoritos no dispositivo.\\
RF-07 & O sistema deve apresentar mensagens de erro quando a API não estiver disponível.\\
\bottomrule
\end{longtable}

\subsection{Requisitos Não Funcionais}

\begin{itemize}
\item RNF-01: O aplicativo deve ser desenvolvido em React Native.
\item RNF-02: O aplicativo deve consumir uma API externa.
\item RNF-03: A interface deve ser simples e responsiva.
\item RNF-04: O sistema deve apresentar feedback de carregamento.
\item RNF-05: O aplicativo deve permitir execução em ambiente Android.
\end{itemize}

\section{Regras de Negócio}

\begin{longtable}{p{2cm}p{8cm}p{2cm}}
\toprule
Código & Descrição & Prioridade\\
\midrule
RN-01 & O usuário deve informar um termo para realizar a busca. & Alta\\
RN-02 & Quando não houver resultados, o sistema deve apresentar uma mensagem. & Alta\\
RN-03 & Uma receita pode ser adicionada ou removida dos favoritos. & Alta\\
RN-04 & Os favoritos devem permanecer salvos no dispositivo. & Média\\
RN-05 & O sistema deve informar o usuário quando ocorrer falha na API. & Alta\\
\bottomrule
\end{longtable}

\section{UML e Modelagem}

\subsection{Diagrama de Caso de Uso}

Os principais casos de uso são:

\begin{itemize}
\item Buscar receita;
\item Visualizar detalhes;
\item Consultar ingredientes;
\item Consultar modo de preparo;
\item Favoritar receita;
\item Consultar favoritos.
\end{itemize}

\subsection{Diagrama de Classes Simplificado}

As principais entidades e módulos são:

\begin{itemize}
\item Usuário: interage com o aplicativo.
\item Receita: contém nome, imagem, categoria, ingredientes e instruções.
\item API: fornece os dados das receitas.
\item Favoritos: realiza a persistência local usando AsyncStorage.
\end{itemize}

\subsection{Diagrama de Sequência Descritivo}

\begin{enumerate}
\item O usuário informa um termo de busca.
\item O aplicativo envia uma requisição HTTP para a TheMealDB.
\item A API retorna uma lista de receitas.
\item O aplicativo apresenta os resultados em cards.
\item O usuário seleciona uma receita.
\item O aplicativo consulta os detalhes e apresenta as informações.
\end{enumerate}

\subsection{Diagrama de Atividades Descritivo}

\begin{enumerate}
\item Início.
\item Informar termo de busca.
\item Validar termo.
\item Consultar API.
\item Verificar resultado.
\item Exibir receitas ou mensagem de erro.
\item Finalizar.
\end{enumerate}

\section{Papéis e Responsabilidades}

Em um projeto individual, o aluno acumula os papéis de Product Owner, desenvolvedor mobile, UX/UI Designer, arquiteto, QA e gerente de projeto.

\begin{longtable}{p{4cm}p{8cm}}
\toprule
Papel & Responsabilidade\\
\midrule
Product Owner & Definir e priorizar os requisitos.\\
Desenvolvedor Mobile & Implementar telas, navegação e integração com a API.\\
UX/UI Designer & Definir uma interface simples e acessível.\\
Arquiteto de Software & Definir tecnologias e organização do projeto.\\
QA/Tester & Executar testes funcionais e de usabilidade.\\
Gerente de Projeto & Organizar cronograma e entregas.\\
\bottomrule
\end{longtable}

\section{Cronograma}

\begin{longtable}{p{2cm}p{5cm}p{5cm}}
\toprule
Semanas & Fase & Entregáveis\\
\midrule
1--2 & Planejamento & Requisitos, UML inicial e cronograma.\\
3--4 & Design & Wireframes e escolha da API.\\
5--8 & Desenvolvimento 1 & Telas, navegação e integração.\\
9--11 & Desenvolvimento 2 & Favoritos, persistência e ajustes.\\
12--13 & Testes & Testes funcionais e avaliação heurística.\\
14 & Build & Geração e teste da APK.\\
15 & Documentação & Revisão do documento e diagramas.\\
16 & Apresentação & Demonstração do aplicativo.\\
\bottomrule
\end{longtable}

\section{Aplicabilidade}

O público-alvo é composto por pessoas que desejam encontrar receitas culinárias de maneira rápida por meio de um dispositivo móvel.

O problema abordado é a dificuldade de localizar receitas e consultar seus ingredientes e instruções em um único local.

A solução é viável porque utiliza uma API pública, bibliotecas consolidadas e armazenamento local para os favoritos.

\section{Arquitetura e Tecnologias}

A aplicação utiliza uma arquitetura simples baseada em componentes e telas.

\begin{itemize}
\item React Native e Expo;
\item Axios para requisições HTTP;
\item TheMealDB como API externa;
\item React Navigation para navegação;
\item FlatList para listagem;
\item AsyncStorage para favoritos.
\end{itemize}

A organização de pastas é:

\begin{verbatim}
src/
  screens/
  services/
  storage/
App.js
\end{verbatim}

\section{Integração com API}

A API utilizada é a TheMealDB.

\begin{longtable}{p{4cm}p{2cm}p{6cm}}
\toprule
Endpoint & Método & Finalidade\\
\midrule
/search.php?s=termo & GET & Buscar receitas por nome.\\
/lookup.php?i=id & GET & Consultar detalhes de uma receita.\\
/filter.php?i=ingrediente & GET & Buscar receitas por ingrediente.\\
/filter.php?c=categoria & GET & Filtrar receitas por categoria.\\
\bottomrule
\end{longtable}

O aplicativo utiliza Axios com async/await e apresenta uma mensagem amigável quando ocorre falha de comunicação ou ausência de resultados.

\section{Usabilidade e UX}

Foram consideradas as seguintes heurísticas de Nielsen:

\begin{longtable}{p{3cm}p{7cm}p{3cm}}
\toprule
Heurística & Aplicação & Evidência\\
\midrule
Visibilidade do status & Indicador de carregamento durante a busca. & Print da tela\\
Correspondência com o mundo real & Uso de nomes, imagens e ingredientes de receitas. & Print da tela\\
Controle do usuário & Permissão para favoritar e remover favoritos. & Print da tela\\
Consistência e padrões & Botões e cards com padrão visual semelhante. & Print da tela\\
Prevenção de erros & Validação de busca vazia e mensagens de erro. & Print da tela\\
Reconhecimento & Exibição visual das receitas em cards. & Print da tela\\
Estética minimalista & Interface com poucos elementos e navegação direta. & Print da tela\\
\bottomrule
\end{longtable}

As capturas de tela devem ser inseridas após a execução real do aplicativo.

\section{Build e Deploy da APK}

A geração da APK deve ser feita no ambiente Android configurado para o projeto.

Procedimento geral:

\begin{enumerate}
\item Instalar as dependências com \texttt{npm install}.
\item Executar o projeto e verificar as telas.
\item Configurar o ambiente de build Android.
\item Gerar a APK de release conforme as orientações da disciplina.
\item Instalar a APK em dispositivo físico ou emulador.
\item Registrar capturas de tela da instalação e execução.
\item Inserir o arquivo APK na pasta \texttt{/release}.
\end{enumerate}

O comando indicado no documento da disciplina é:

\begin{verbatim}
cd android
./gradlew clean
./gradlew assembleRelease
\end{verbatim}

O comando deve ser utilizado quando o projeto estiver configurado com a estrutura Android necessária para o build de release.

\section{Conclusão e Lições Aprendidas}

O projeto Receita Fácil apresenta uma solução simples para consulta de receitas culinárias. A implementação utiliza React Native, integração com uma API externa e armazenamento local para favoritos.

Durante o desenvolvimento, os principais aprendizados envolvem consumo de APIs, organização de componentes, navegação entre telas, persistência de dados e aplicação de princípios de usabilidade.

Como melhorias futuras, podem ser adicionados filtros mais completos, geração automática de lista de compras, autenticação de usuários e sincronização dos favoritos em nuvem.

\section{Referências}

\begin{itemize}
\item TheMealDB. Disponível em: \url{https://www.themealdb.com/api.php}.
\item React Native. Documentação oficial. Disponível em: \url{https://reactnative.dev/}.
\item React Navigation. Documentação oficial. Disponível em: \url{https://reactnavigation.org/}.
\item Nielsen Norman Group. 10 Usability Heuristics for User Interface Design.
\end{itemize}

\end{document}
